\chapter{متجهات الحالة، والتطور الوحدوي، وتناظر الإزاحة}
\providecommand{\ket}[1]{\left\lvert#1\right\rangle}
\providecommand{\bra}[1]{\left\langle#1\right\rvert}

بعد انقطاع دام ربعين دراسيين، تكون مراجعة سريعة لميكانيكا الكم مفيدة. ولا يمكن أن تستغرق الربع الدراسي كله. والموضوع القادم هو تطبيق ميكانيكا الكم على المسائل التي طُوِّرت أصلًا من أجلها: الذرات، والإلكترونات والجسيمات المرتبطة بها، والفوتونات.

تأتي الذرات أولًا. كيف يمكن أن تأتي الذرات قبل الإلكترونات، مع أن الإلكترونات أجزاء من الذرات؟ لا نحتاج في البداية من ميكانيكا الإلكترون إلا إلى القدر الكافي لكتابة المعادلات الأساسية لذرة الهيدروجين. وبعد ذلك يمكن أن تتبعها دراسة أكثر تفصيلًا للجسيمات الشبيهة بالإلكترون. وفي جميع هذه التطبيقات سيكون التناظر أداة متكررة: كيف تُمثَّل التناظرات في ميكانيكا الكم، وما الذي تقتضيه بشأن نظام ما.

\section{الحالات، والمرصودات، والتعامد}

في الميكانيكا الكلاسيكية، تُمثَّل الحالة بنقطة في فضاء الطور. وفي ميكانيكا الكم، تُمثَّل بمتجه في فضاء متجهي مركب. وهذا متجه مجرد، لا سهم في الفضاء الثلاثي الأبعاد العادي. ويمكن جمع متجهات الحالة وضربها في أعداد مركبة، ولا يلزم أن يكون الفضاء الذي يحتويها ذا أي بعد منتهٍ معين.

يكتب ترميز ديراك متجه الحالة على صورة كيت،
\begin{equation}
\ket{\psi},
\end{equation}
ويربط به برا،
\begin{equation}
\bra{\psi}.
\end{equation}
وترتبط البرا والكيت بأخذ المرافق الهرميتي. وضم برا إلى كيت يعطي جداءً داخليًا،
\begin{equation}
\langle\phi|\psi\rangle,
\end{equation}
ويكون هذا عمومًا عددًا مركبًا.

المرصود هو شيء يمكن قياسه: الموضع، أو الزخم، أو الطاقة، أو أي كمية أخرى قابلة للقياس. وفي ميكانيكا الكم يُمثَّل المرصود بمؤثر خطي هرميتي \(A\)،
\begin{equation}
A^\dagger=A.
\end{equation}
والهرميتية هي النظير على مستوى المؤثرات لكون الكمية حقيقية. وعلى وجه الخصوص، فإن القيم الذاتية لمؤثر هرميتي تكون حقيقية.

متجه ذاتي لـ \(A\) يحقق
\begin{equation}
A\ket{\alpha}=\alpha\ket{\alpha}.
\end{equation}
والقيم الذاتية الممكنة هي النتائج الممكنة لقياس \(A\). فإذا كان مرصود ما لا يمكن أن يعطي إلا \(1\)، أو \(3\)، أو \(7\)، فإن هذه الأعداد هي قيمه الذاتية. أما متجهاته الذاتية فهي الحالات التي تكون فيها النتيجة محددة لا إحصائية. وفي الحالة \(\ket{\alpha}\)، فإن قياس \(A\) يعطي \(\alpha\) على وجه اليقين.

يعبِّر الجداء الداخلي عن علاقة منطقية مهمة بين الحالات. افترض وجود قياس يميز \(\ket{\phi}\) عن \(\ket{\psi}\) على وجه اليقين. عندئذ تكون الحالتان متعامدتين:
\begin{equation}
\langle\phi|\psi\rangle=0.
\end{equation}
ويعني التعامد أن الحالتين يمكن تمييزهما فيزيائيًا بلا لبس. ولحالة مطبَّعة،
\begin{equation}
\langle\psi|\psi\rangle=1.
\end{equation}

\section{جسيم في تمثيل الموضع}

لننظر أولًا في جسيم يتحرك على مستقيم. موضعه \(x\) مرصود. ولكل موضع ممكن \(x_0\)، توجد حالة \(\ket{x_0}\) يكون فيها قياس الموضع معطيًا يقينًا \(x_0\).

الحالات المرتبطة بمواضع مختلفة قابلة للتمييز، ومن ثم فهي متعامدة:
\begin{equation}
x\neq x'
\quad\Longrightarrow\quad
\langle x'|x\rangle=0.
\end{equation}
ولحالة اعتباطية \(\ket{\psi}\)، يعرِّف تراكبها مع حالة موضع الدالة الموجية:
\begin{equation}
\psi(x)=\langle x|\psi\rangle.
\end{equation}
وهذا هو تمثيل الموضع لمتجه الحالة.

الدالة الموجية ليست هي نفسها احتمالًا. يمكن أن تكون موجبة، أو سالبة، أو تخيلية، أو مركبة عمومًا. وكثافة احتمال الموضع هي
\begin{equation}
P(x)=\psi^*(x)\psi(x)=|\psi(x)|^2.
\end{equation}
وهذا الجداء غير سالب، كما يجب لكثافة الاحتمال. وللدالة الموجية المطبَّعة،
\begin{equation}
\int_{-\infty}^{\infty}dx\,|\psi(x)|^2=1.
\end{equation}

في تمثيل الموضع، يؤثر مؤثر الموضع بالضرب:
\begin{equation}
(\hat x\psi)(x)=x\psi(x).
\end{equation}
ودواله الذاتية دوال مثالية متمركزة في نقطة واحدة. وتوصَف صوريًا بأنها دوال دلتا لديراك: قمم ضيقة جدًا متمركزة عند القيمة الذاتية للموضع. وليس تعريفها التوزيعي الدقيق مطلوبًا في هذه المراجعة الحالية.

في ثلاثة أبعاد، للموضع ثلاث مركبات، \(x\) و\(y\) و\(z\)، وتمثل حالة الموضع نقطة في الفضاء الثلاثي الأبعاد. وتكون حالتان من هذا النوع متعامدتين كلما اختلفت النقطتان في إحداثي واحد على الأقل. وسيُحتفَظ عادةً بالترميز أحادي البعد، مع فهم التعميم الثلاثي الأبعاد المناظر.

\section{الزخم وحالات الزخم المحدد}

يقابل كل مركبة من مركبات الموضع مركبة مناظرة من الزخم. وفي تمثيل الموضع يكون مؤثر الزخم على طول محور \(x\) هو
\begin{equation}
\hat p_x=-i\hbar\frac{d}{dx}.
\end{equation}
وهكذا يمكن أيضًا النظر إلى مؤثر هرميتي مجرد على نحو ملموس بوصفه عملية على الدوال الموجية.

الحالة ذات الزخم المحدد \(p_0\) تحقق معادلة القيم الذاتية
\begin{equation}
\hat p_x\psi_{p_0}(x)=p_0\psi_{p_0}(x).
\end{equation}
وبالتعويض عن الصورة التفاضلية لـ \(\hat p_x\) نحصل على
\begin{equation}
-i\hbar\frac{d\psi_{p_0}(x)}{dx}
=p_0\psi_{p_0}(x).
\end{equation}
ومشتقة الدالة الموجية تتناسب مع الدالة الموجية نفسها، لذا يكون الحل ذا صورة أسية:
\begin{equation}
\psi_{p_0}(x)\propto
\exp\left(\frac{ip_0x}{\hbar}\right).
\end{equation}

\begin{classroomqa}
\lecturetimestamp{00:24:26}
\audiencequestion{هل تصدق هذه المعادلة التفاضلية على أي دالة موجية؟}
\lecturerresponse{لا. إنها معادلة القيمة الذاتية لحالة ذات زخم محدد. ومعظم الدوال الموجية لا تحققها؛ فهي تنتقي دوال الزخم الذاتية.}
\end{classroomqa}

التمييز بين حالات الموضع الذاتية وحالات الزخم الذاتية حاد. فدالة الموضع الذاتية متمركزة في نقطة واحدة. أما دالة الزخم الذاتية فتحقق
\begin{equation}
\left|\exp\left(\frac{ip_0x}{\hbar}\right)\right|^2=1.
\end{equation}
ولذلك يكون معيارها ثابتًا على امتداد المستقيم كله. وفي هذا المعنى المثالي، يعني الزخم المحدد لايقينًا كاملًا في الموضع؛ وكذلك يعني الموضع المحدد لايقينًا كاملًا في الزخم. وهذه هي الصورة القصوى لعلاقة الارتياب في الموضع--الزخم.

\begin{classroomqa}
\lecturetimestamp{00:26:36}
\audiencequestion{إذا كان معلومًا أن الزخم هو \(p_0\)، فلماذا لا يزال هناك توزيع احتمالي؟}
\lecturerresponse{الكمية المعروضة تتعلق بقياس الموضع، لا بقياس آخر للزخم. فالزخم والموضع مرصودان مختلفان. وبالنسبة إلى حالة الزخم الذاتية المثالية، فإن سعة الموضع لها معيار ثابت على امتداد المستقيم كله. أما سعات الزخم فتنتمي إلى تمثيل الزخم، وترتبط بالدالة الموجية في الموضع بتحويل فورييه.}
\end{classroomqa}

\begin{classroomqa}
\lecturetimestamp{00:28:02}
\audiencequestion{كيف تبدو كثافة احتمال الموضع عندما تكون الدالة الموجية دالة دلتا لديراك؟}
\lecturerresponse{إنها متمركزة عند الموضع المحدد. وتعبر القمة المثالية عن اليقين في الموضع.}
\end{classroomqa}

\editorialnote{حالات القيم الذاتية الدقيقة للموضع والزخم هي حالات معمَّمة. ودالة دلتا لديراك وموجة مستوية ليستا دالتين موجيتين عاديتين قابلتين للتكامل تربيعيًا ومطبعتين على المستقيم كله، ومربع دالة دلتا بمعناه الحرفي ليس توزيعًا عاديًا. وتتعلق العبارات أعلاه بتمركزهما المثالي وثبات معيارهما.}

يمكن النظر إلى الدالة الموجية صوريًا على أنها متجه له مركبة واحدة لكل قيمة من قيم \(x\). لكنها، بخلاف متجه عمودي منتهٍ عادي، تملك عددًا لا نهائيًا متصلًا من المركبات.

\begin{classroomqa}
\lecturetimestamp{00:28:53}
\audiencequestion{هل تُمثَّل متجهات الحالة بمتجهات عمودية؟}
\lecturerresponse{يمكن تمثيلها على هذا النحو، لكن كتابة دالة موجية موضعية على شكل عمود عادي غير مريحة لأن \(x\) يأخذ عددًا متصلًا من القيم. ويمكن النظر إليها صوريًا بوصفها متجهًا عموديًا متصلًا، مع اعتبار البرا الموافق متجه صف متصلًا.}
\end{classroomqa}

\section{التطور الزمني والوحدوية}

كيف تتغير الحالة مع الزمن؟ التطور الزمني هو تحويل على فضاء الحالات. ولْيَدُلَّ \(U(t)\) على التطور خلال زمن منقضٍ مقداره \(t\). إذًا
\begin{equation}
U(t)\ket{\psi(t_1)}=\ket{\psi(t_1+t)}.
\end{equation}
وعلى وجه الخصوص،
\begin{equation}
\ket{\psi(t)}=U(t)\ket{\psi(0)},
\qquad
U(0)=\mathbf 1.
\end{equation}

ينتج الشرط على \(U(t)\) من حفظ المعلومات. فإذا كانت حالتان ابتدائيتان قابلتين للتمييز، فلا بد ألا يمحو التطور الزمني ذلك التمييز. ويجب أن تظل الحالات المتعامدة متعامدة. وعلى نحو أعم، يُحفَظ الجداء الداخلي بين أي حالتين:
\begin{equation}
\langle U(t)\phi|U(t)\psi\rangle
=\langle\phi|\psi\rangle.
\end{equation}
وبما أن
\begin{equation}
\langle U(t)\phi|U(t)\psi\rangle
=\langle\phi|U^\dagger(t)U(t)|\psi\rangle,
\end{equation}
فإن حفظ جميع الجداءات الداخلية يقتضي
\begin{equation}
U^\dagger(t)U(t)=\mathbf 1.
\end{equation}
والمؤثر الذي يحقق هذه المعادلة يكون وحدويًا.

\begin{classroomqa}
\lecturetimestamp{00:36:36}
\audiencequestion{هل يُعَدُّ حفظ كل جداء داخلي عبارةً أقوى من مجرد حفظ التعامد؟}
\lecturerresponse{بالنسبة إلى تحويلات الحالات الخطية التي نناقشها، فإن حفظ التعامد يفضي إلى حفظ الجداءات الداخلية. وتقدم الدورانات العادية القياس المألوف: فهي تحفظ التعامد والضرب النقطي.}
\end{classroomqa}

\editorialnote{على وجه التدقيق، فإن حفظ التعامد وحده يسمح بإعادة تحجيم كلية غير صفرية. أما حفظ التطبيع، وهو مطلوب أصلًا للتطور الزمني الفيزيائي، فيزيل تلك الحرية ويعطي حفظ الجداءات الداخلية.}

لا ينبغي الخلط بين المؤثرات الوحدوية والمرصودات الهرميتية. فالموضع والزخم هرميتيان، لكنهما ليسا عمومًا وحدويين. والقيم الذاتية لمؤثر وحدوي هي أطوار، أي أعداد من الصورة \(e^{i\theta}\) تقع على دائرة الوحدة.

\begin{classroomqa}
\lecturetimestamp{00:38:12}
\audiencequestion{ما مثال على تحول غير وحدوي؟}
\lecturerresponse{مؤثرا الموضع والزخم مثالان على ذلك: فهما مؤثران هيرميتيان يمثّلان كميتين قابلتين للرصد، لكنهما ليسا وحدويين. للمؤثر الوحدوي قيم ذاتية معيارها واحد، لذا تكون قيمه الذاتية عوامل طور.}
\end{classroomqa}

\section{الهاملتوني ومعادلة Schr\"odinger}

يُفترض أن التطور الزمني متصل. ولذلك، على فاصل قصير جدًا \(\epsilon\)، لا بد أن يكون مؤثر التطور قريبًا من الهوية:
\begin{equation}
U(\epsilon)=\mathbf 1+\epsilon G+\mathcal O(\epsilon^2).
\end{equation}
وباستخدام الوحدوية،
\begin{align}
U^\dagger(\epsilon)U(\epsilon)
&=
\left(\mathbf 1+\epsilon G^\dagger\right)
\left(\mathbf 1+\epsilon G\right)
+\mathcal O(\epsilon^2)
\\
&=
\mathbf 1+\epsilon(G^\dagger+G)
+\mathcal O(\epsilon^2).
\end{align}
وإلى الرتبة الأولى،
\begin{equation}
G^\dagger=-G.
\end{equation}
إذًا \(G\) مضادّ هيرميتي. ويمكن كتابته على الصورة
\begin{equation}
G=-\frac{i}{\hbar}H,
\end{equation}
حيث إن \(H\) هيرميتي:
\begin{equation}
H^\dagger=H.
\end{equation}
ولذلك يكون مؤثر التطور الزمني المتناهي في الصغر
\begin{equation}
U(\epsilon)
=
\mathbf 1-\frac{i\epsilon}{\hbar}H
+\mathcal O(\epsilon^2).
\end{equation}
والمؤثر الهيرميتي \(H\) هو الهاملتوني.

\begin{classroomqa}
\lecturetimestamp{00:45:48}
\audiencequestion{ما الذي يجيز إهمال الحد \(\epsilon^2\)؟}
\lecturerresponse{يُجرى الحساب باتساق حتى الرتبة الأولى في \(\epsilon\). أما الحساب من الرتبة الثانية فسيتطلب جميع حدود الرتبة الثانية، بما فيها الحد المتناسب مع \(-\epsilon^2H^2/(2\hbar^2)\). يقع التقريب في التعبير المقطوع لـ \(U(\epsilon)\)، لا في هيرميتية \(H\).}
\end{classroomqa}

يمكن بناء فاصل زمني منتهٍ من كثير من الفواصل الصغيرة. اقسم \(t\) إلى \(N\) خطوات طول كل منها \(t/N\). ويعطي تكرار التحول المتناهي في الصغر
\begin{equation}
U(t)
=
\lim_{N\to\infty}
\left(
\mathbf 1-\frac{itH}{N\hbar}
\right)^N
=
\exp\left(-\frac{iHt}{\hbar}\right).
\end{equation}
والأسي هو الصيغة المنتهية للتحول المتناهي في الصغر.

لنطوّر الآن حالة من \(t\) إلى \(t+\epsilon\):
\begin{align}
\ket{\psi(t+\epsilon)}
&=
U(\epsilon)\ket{\psi(t)}
\\
&=
\left(
\mathbf 1-\frac{i\epsilon}{\hbar}H
\right)\ket{\psi(t)}
+\mathcal O(\epsilon^2).
\end{align}
وبطرح \(\ket{\psi(t)}\)، ثم القسمة على \(\epsilon\)، ثم أخذ النهاية، نحصل على
\begin{equation}
\frac{d}{dt}\ket{\psi(t)}
=
-\frac{i}{\hbar}H\ket{\psi(t)}.
\end{equation}
وبصورة مكافئة،
\begin{equation}
i\hbar\frac{d}{dt}\ket{\psi(t)}
=
H\ket{\psi(t)}.
\end{equation}
هذه هي معادلة Schr\"odinger المعتمدة على الزمن. وفي تمثيل الموضع تصبح
\begin{equation}
i\hbar\frac{\partial}{\partial t}\psi(x,t)
=
H\psi(x,t),
\end{equation}
حيث يعمل \(H\) بوصفه مؤثرًا على الدالة الموجية.

معادلة Schr\"odinger المستقلة عن الزمن هي معادلة القيم الذاتية للهاملتوني:
\begin{equation}
H\ket{E}=E\ket{E}.
\end{equation}
والقيم الذاتية \(E\) هي الطاقات الممكن قياسها، أي مستويات الطاقة. ولذا يحدد الهاملتوني نفسه كلًّا من الطاقات المسموح بها وتطور الحالة.

إذا كانت الحالة الابتدائية حالةً ذاتية للطاقة، فإن تطورها هو
\begin{equation}
\ket{E;t}
=
\exp\left(-\frac{iEt}{\hbar}\right)\ket{E}.
\end{equation}
فيبقى متجه الطاقة الذاتي نفسه ولا يكتسب إلا طورًا.\footnote{تتحدد الإشارة السالبة من الاشتقاق السابق لـ \(U(t)=\exp(-iHt/\hbar)\). وفي صياغة منطوقة لاحقة للطور تظهر الإشارة المعاكسة، وهو ما لا يتسق مع ذلك الاشتقاق.}

\begin{classroomqa}
\lecturetimestamp{01:01:05}
\audiencequestion{كيف يتطور متجه ذاتي للطاقة وفق معادلة Schr\"odinger المعتمدة على الزمن؟}
\lecturerresponse{احتفظ بالكيت. بما أن \(H\ket{E}=E\ket{E}\)، فإن التطور الزمني يضرب متجه الطاقة الذاتي بطور تحدده طاقته: \(\exp(-iEt/\hbar)\).}
\end{classroomqa}

\section{ما الذي يُعَدّ تناظرًا؟}

يمثل التطور الزمني أحد التحولات الوحدوية على فضاء الحالات. وتشكل تحولات التناظر فئة مهمة أخرى.

يحتل تناظر الدوران مكانة مركزية في الفيزياء الذرية. ففي ذرة الهيدروجين يتحرك الإلكترون في جهد كولوم المركزي. إن تدوير جهاز الإحداثيات لا يغير المعادلات، كما أن تدوير الذرة نفسها فيزيائيًا يتركها أيضًا ذرة هيدروجين تحكمها المعادلات نفسها. قد تتغير الحالة تحت الدوران، لكن القانون الديناميكي لا يتغير.

يعطي الانسحاب مثالًا آخر. ينبغي أن يسلك نظام معزول في فضاء خالٍ السلوك نفسه بعد أن تُزاح كل أجزائه بالمقدار نفسه. وهذا القيد مهم. فإذا وُجد فرن ونُقل النظام وحده ليصبح أقرب إليه، فقد يسلك النظام سلوكًا مختلفًا؛ فالفرن يكسر تناظر الانسحاب ذاك. أما إذا أُزيح كل من النظام والفرن معًا، فإن العلاقة الأصلية بينهما تُستعاد.

ليست كل عملية تناظرًا. فضغط بلورة من الملح الصخري بمعامل اثنين يغير سلوكها عمومًا. وهذا الضغط تحول محدد جيدًا، لكنه ليس تناظرًا على العموم.

وعليه فالتناظر عملية على الحالات تحفظ صورة الديناميكيات. ولنرمز إليه بـ \(V\). ويجب أن يحفظ قابلية التمييز: فلا بد أن تبقى حالتان متعامدتان متعامدتين بعد التحول. وبناءً على ذلك،
\begin{equation}
V^\dagger V=\mathbf 1,
\end{equation}
ومن ثم فإن \(V\) وحدوي.

\section{التناظرات تتبادل مع التطور الزمني}

افترض أن حالة ابتدائية \(\ket{\psi_1}\) تتطور إلى \(\ket{\psi_2}\):
\begin{equation}
\ket{\psi_2}=U(t)\ket{\psi_1}.
\end{equation}
طبّق تناظرًا مقترحًا \(V\) على كلتا الحالتين:
\begin{equation}
\ket{\psi_1'}=V\ket{\psi_1},
\qquad
\ket{\psi_2'}=V\ket{\psi_2}.
\end{equation}
إذا كان \(V\) تناظرًا للديناميكيات، فلا بد أن تتطور الحالة الابتدائية المحوَّلة إلى الحالة النهائية المحوَّلة:
\begin{equation}
\ket{\psi_2'}=U(t)\ket{\psi_1'}.
\end{equation}
وبالتعويض عن التعريفات نحصل على
\begin{equation}
VU(t)\ket{\psi_1}
=
U(t)V\ket{\psi_1}.
\end{equation}
ويجب أن يصح هذا لكل حالة ابتدائية. لذلك
\begin{equation}
VU(t)=U(t)V.
\end{equation}

يمكن عرض المسارين المتكافئين في مخطط تبادلي.\footnote{يعيد هذا المخطط بناء حجة المسارين المطوَّرة بين 01:10:30 و01:16:34.}
\begin{figure}
\centering
\begin{tikzpicture}[>=Stealth]
\node (a) at (0,0) {$\ket{\psi_1}$};
\node (b) at (4.2,0) {$\ket{\psi_2}=U(t)\ket{\psi_1}$};
\node (c) at (0,-2.0) {$V\ket{\psi_1}$};
\node (d) at (4.2,-2.0) {$V\ket{\psi_2}=U(t)V\ket{\psi_1}$};

\draw[->] (a) -- node[above] {$U(t)$} (b);
\draw[->] (c) -- node[below] {$U(t)$} (d);
\draw[->] (a) -- node[left] {$V$} (c);
\draw[->] (b) -- node[right] {$V$} (d);
\end{tikzpicture}
\caption{يعطي تطبيق تناظر والتطور الزمني الحالة النهائية نفسها في أيٍّ من الترتيبين.}
\end{figure}

\begin{classroomqa}
\lecturetimestamp{01:16:38}
\audiencequestion{هل يجب أن يتبادل التناظر مع التطور الزمني لجميع الأزمنة؟}
\lecturerresponse{نعم. يجب أن يحقق تناظر الديناميكيات شرط التبادل مع التطور الزمني الكامل.}
\end{classroomqa}

\begin{classroomqa}
\lecturetimestamp{01:17:13}
\audiencequestion{هل يعني هذا أنه يمكن إجراء دوران ثم الانتظار، أو الانتظار ثم إجراء دوران؟}
\lecturerresponse{بالضبط. إذا كان الدوران تناظرًا، فإن كلا الترتيبين يعطي الحالة النهائية المحوَّلة نفسها.}
\end{classroomqa}

على فاصل زمني صغير،
\begin{equation}
U(\epsilon)
=
\mathbf 1-\frac{i\epsilon}{\hbar}H
+\mathcal O(\epsilon^2).
\end{equation}
ويصبح شرط التبادل
\begin{equation}
V\left(
\mathbf 1-\frac{i\epsilon}{\hbar}H
\right)
=
\left(
\mathbf 1-\frac{i\epsilon}{\hbar}H
\right)V
+\mathcal O(\epsilon^2).
\end{equation}
وتتلاشى حدود الهوية، ويبقى
\begin{equation}
VH=HV,
\end{equation}
أو
\begin{equation}
[V,H]=0.
\end{equation}
إذًا تكون العملية الوحدوية تناظرًا للديناميكيات إذا وفقط إذا كانت تتبادل مع الهاملتوني.

\begin{classroomqa}
\lecturetimestamp{01:19:29}
\audiencequestion{ما الذي يقتضي أن يكون مؤثر التناظر \(V\) وحدويًا؟}
\lecturerresponse{يجب على التناظر أن يحفظ الفروق المنطقية بين الحالات. فلا ينبغي له أن يحوّل حالتين متنافيتين ومتعامدتين إلى حالتين لا تعودان قابلتين للتمييز. وحفظ تلك العلاقات يتطلب تحولًا وحدويًا.}
\end{classroomqa}

المؤثر الذي لا يملك اعتمادًا زمنيًا صريحًا ويتبادل مع \(H\) يكون محفوظًا. وهذا يربط التناظر بالانحفاظ. وفي الميكانيكا الكلاسيكية تُعبَّر هذه العلاقة بمبرهنة Noether. وفي الصيغة الكمومية الحالية، المعيار الأساسي هو المبدّل مع الهاملتوني.

\section{التناظرات المتقطعة والمستمرة}

يمكن أن تكون التناظرات متقطعة أو مستمرة. فالانعكاس المرآتي متقطع: فهو يبدّل اليسار باليمين، من غير وجود سلسلة من انعكاسات متناهية الصغر بينهما. ويُعدّ تبديل جسيمين متماثلين عملية متقطعة أخرى. كما أن قلب جسم ذي وجهين غير معلَّمين يعطي النوع نفسه من الأمثلة.

أما الدورانات فمستمرة. فدوران بمقدار \(90^\circ\) يمكن بناؤه من كثير من الدورانات الصغيرة حول المحور نفسه. وكذلك الانسحابات المكانية مستمرة للسبب نفسه: إذ يمكن بناء إزاحة منتهية من كثير من الإزاحات الصغيرة جدًا.

إن إمكان تجزئة التحول إلى تحولات صغيرة يمنح التناظرات المستمرة صورة خاصة. وقرب الهوية، نكتب
\begin{equation}
V(\epsilon)
=
\mathbf 1-i\epsilon G+\mathcal O(\epsilon^2),
\end{equation}
حيث يحدد \(\epsilon\) التحول. وتقتضي الوحدوية
\begin{equation}
G^\dagger=G.
\end{equation}
ويُسمّى المؤثر الهيرميتي \(G\) مولِّد التناظر.

إذا كان التحول تناظرًا، فإن
\begin{equation}
[H,V(\epsilon)]=0.
\end{equation}
ويعطي التوسيع إلى الرتبة الأولى
\begin{equation}
[H,G]=0.
\end{equation}
ولذلك تُولَّد التناظرات المستمرة بمؤثرات هيرميتية تتبادل مع الهاملتوني. وهذه المولِّدات كميات محفوظة.

لكلمة ``مولَّدة'' هنا معنى ملموس. فالتحول المنتهي يُبنى بتطبيق التحول المتناهي في الصغر مرارًا، تمامًا كما يُبنى التطور الزمني المنتهي من كثير من الخطوات الزمنية القصيرة.

\section{تناظر الانسحاب والزخم}

اعتبر انسحابًا فعّالًا لدالة موجية إلى اليمين بمسافة \(\epsilon\). فإذا كانت الدالة الموجية الأصلية تتمركز عند \(x_0\)، فينبغي أن تتمركز الدالة الموجية المنساحة عند \(x_0+\epsilon\). ولذلك تكون قيمتها عند الإحداثي \(x\) هي قيمة الدالة الموجية الأصلية المحسوبة عند \(x-\epsilon\):
\begin{equation}
(T(\epsilon)\psi)(x)=\psi(x-\epsilon).
\end{equation}
ولإزاحة صغيرة،
\begin{equation}
\psi(x-\epsilon)
=
\psi(x)-\epsilon\frac{d\psi(x)}{dx}
+\mathcal O(\epsilon^2).
\end{equation}
ومن ثم
\begin{equation}
T(\epsilon)
=
\mathbf 1-\epsilon\frac{d}{dx}
+\mathcal O(\epsilon^2).
\end{equation}
وباستخدام
\begin{equation}
\hat p_x=-i\hbar\frac{d}{dx},
\qquad
\frac{d}{dx}=\frac{i}{\hbar}\hat p_x,
\end{equation}
يصبح الانسحاب المتناهي في الصغر
\begin{equation}
T(\epsilon)
=
\mathbf 1-\frac{i\epsilon}{\hbar}\hat p_x
+\mathcal O(\epsilon^2).
\end{equation}
والانسحاب المنتهي الموافق هو
\begin{equation}
T(\epsilon)
=
\exp\left(-\frac{i\epsilon}{\hbar}\hat p_x\right).
\end{equation}
ومع الاصطلاح \(V(\epsilon)=\mathbf 1-i\epsilon G+\cdots\)، يعطي وسيط الانسحاب ذي الأبعاد
\begin{equation}
G_{\mathrm{trans}}=\frac{\hat p_x}{\hbar}.
\end{equation}
وبصورة مكافئة، يُسمَّى الزخم مولِّد الانسحابات المكانية عندما يُظهَر عامل \(\hbar\) في الأس.

\begin{classroomqa}
\lecturetimestamp{01:34:08}
\audiencequestion{هل فُقدت إشارة سالبة في اشتقاق الانسحاب؟}
\lecturerresponse{المسألة هي الاتجاه الذي يُسنَد إلى الدالة الموجية المنساحة. فالانسحاب الفعّال إلى اليمين يُمثَّل بـ \(\psi(x-\epsilon)\)، لا بـ \(\psi(x+\epsilon)\). ومع هذا التصحيح، \(T(\epsilon)=\mathbf 1-i\epsilon\hat p_x/\hbar+\cdots\).}
\end{classroomqa}

إنَّ ما إذا كان الزخم محفوظًا أم لا يعتمد على الهاملتوني. بالنسبة إلى جسيم حر يتحرك في بعد واحد،
\begin{equation}
H=\frac{\hat p_x^2}{2m}.
\end{equation}
وبما أن كل مؤثر يتبادل مع أي دالة له،
\begin{equation}
[H,\hat p_x]=0.
\end{equation}
فإن هاملتوني الجسيم الحر لا يتغير تحت الانسحابات، ومن ثم يكون \(\hat p_x\) محفوظًا. وفيزيائيًا، يسلك الجسيم السلوك نفسه أيًّا كان الموضع الابتدائي المختار له. وإن إخضاع الشرط الابتدائي لانسحاب لا يفعل سوى إزاحة الحركة بأكملها.

في ثلاثة أبعاد،
\begin{equation}
H
=
\frac{\hat p_x^2+\hat p_y^2+\hat p_z^2}{2m}.
\end{equation}
كل مركبة من مركبات الزخم تتبادل مع \(H\):
\begin{equation}
[H,\hat p_x]=[H,\hat p_y]=[H,\hat p_z]=0.
\end{equation}
يوجد تناظر انسحابي مستقل في كل اتجاه مكاني. والزخم هو في آن واحد الكمية المحفوظة والمولِّد المرتبط بتلك الانسحابات.

الحالة التالية هي التناظر الدوراني. والدورانات أكثر تعقيدًا لأن الدورانات حول محاور مختلفة لا تتبادل عمومًا: فدوران حول \(x\) يتبعه دوران حول \(y\) ليس هو نفسه في الترتيب العكسي. وهذا يؤدي إلى مولدات تناظر غير متبادلة وإلى البنية الزمرية للدورانات.
